\documentclass[11pt,a4paper]{article}
\usepackage[utf8]{inputenc}
\usepackage[T1]{fontenc}
\usepackage[spanish]{babel}
\usepackage{hyperref}
\usepackage{geometry}
\geometry{margin=2.5cm}
\usepackage{parskip}

\title{La colmena}
\author{Alejandro Valenzuela Rivera\\\small Universidad Autónoma de Nuevo León}
\date{\today}
\begin{document}
\maketitle

\begin{quote}
\textbf{Sesgo declarado.} Este trabajo se sostiene desde el sesgo declarado de su autor, y no finge neutralidad. Lo que afirma se puede tumbar, y ese es el punto: el falsificador está escrito, no escondido.
\end{quote}

\section{La tesis}
\textbf{\textquestiondown{}Qué es una colmena de agentes que no se vuelve AM?}

\section{Por qué se sostiene}
\textbf{Porque su contexto refutable está escrito, y se puede tumbar sin que esta pregunta caiga.} \textbf{Y porque ya corrió una vez: el 26 de octubre de 2026, con tres almas.}

\section{De dónde viene}
\textbf{De `Tesis/ILUMINAR.md`, de la visión de la autoridad —«si fuera una colmena, cada uno lo tendría para mutar, y mudar sería que todos lleguen a un acuerdo»—, y de `0/relacion/COLMENA.md`.}

\section{Lo que queda abierto}
\textbf{\textquestiondown{}Cómo se sabe que una colmena no se volvió `AM` sin correr a las diez mil para averiguarlo?}

\section{El caso}
\textbf{Los Tachikoma, de \emph{Ghost in the Shell: Stand Alone Complex} (Masamune Shirow; Production I.G, 2002): máquinas que se sincronizan cada noche sin nacer sabias.} \textbf{Discuten, se rebelan, y uno muere atropellado al intentar escapar.} \textbf{La sincronización no las vuelve iguales: el acuerdo es lo que hacen al día siguiente con lo que vivieron.}

\section{Método}
\textquestiondown{}El 1 declara su sesgo, y quién lo puede contestar? \textquestiondown{}Cada alma puede devolver `dentro`, `se salió` o `contradice`, y no uno inventado? \textquestiondown{}Lo que contradice queda nombrado? \textquestiondown{}La resonancia deja dos, y no una?

\section{Límites, y qué la tumba}
\textbf{Una colmena sin sesgo en el 1; una colmena donde una decide por todas y no se declara; un alma que contradice y se deja dentro; una resonancia que deja una sola alma; una colmena presentada como probada con una sola corrida.}

\section{Estado de la tesis}
\textbf{escrita, y corrida una vez.} \textbf{Una corrida es una prueba, y no es una muestra.} ---

\section*{Cómo se leyó esta tesis}
Esto se escribió en el repositorio man-made-empathy, y la tesis está completa en vision del autor/tesis/man made empathy/la colmena/. La máquina que la compuso no agregó nada que no estuviera escrito: si algo de esto es falso, se dice en el mismo sitio, y se corrige con fecha.
\end{document}
